\documentclass{beamer}

%% Tema y configuraciones
\usetheme{Boadilla}
\usecolortheme{default}
\setbeamertemplate{navigation symbols}{}
\usepackage{xcolor}
\usepackage{tikz}
\usetikzlibrary{arrows.meta, positioning}
\hypersetup{
    colorlinks=true,
    linkcolor=.,      % color de enlaces internos (índices, refs)
    urlcolor=blue,    % color de URLs
    citecolor=.,      % color de referencias bibliográficas
}

%% Helpers (macros y notación)
\newcommand{\E}{\mathbb{E}}
\newcommand{\Var}{\mathrm{Var}}
\newcommand{\Cov}{\mathrm{Cov}}
\newcommand{\1}{\mathbb{I}}
\newcommand{\MSE}{\mathrm{MSE}}
\DeclareMathOperator*{\argmin}{arg\,min}

%% Encabezado
\title[BDML]{Big Data y Machine Learning \\ para Economía Aplicada}
\subtitle{Week 01: Predecir, no explicar --- el marco del aprendizaje estadístico}
\author{Eduard F. Martinez Gonzalez, Ph.D.}
\institute[]{Departamento de Economía, Universidad Icesi}
\date{\today}

%%=================================================
%% Begin document
\begin{document}

%%=================================================
%% Primer slide
\begin{frame}
\titlepage
\end{frame}

%%#################################################
%%#################################################
%% PARTE 1: ACERCA DEL CURSO
%%#################################################
%%#################################################

%%=================================================
\begin{frame}{Acerca de este curso}
\small
\begin{enumerate}\itemsep3pt
\item \textbf{Al finalizar el curso podrás:}
\begin{itemize}\itemsep2pt
  \item Predecir y clasificar variables económicas con métodos de ML y evaluar el desempeño \emph{fuera de muestra} con honestidad.
  \item Construir el flujo de trabajo completo en R: partición, validación, ajuste, evaluación e interpretación.
  \item Abrir un modelo flexible (importancia de variables, dependencia parcial, SHAP) y saber qué dice y qué no dice.
  \item Leer críticamente y presentar investigación económica que usa ML.
\end{itemize}
\item \textbf{Texto guía:} James, Witten, Hastie \& Tibshirani, \emph{An Introduction to Statistical Learning} (ISL, 2.\ ed.). Seguimos su progresión.
\item \textbf{Metodología:} en cada método, el problema que resuelve $\rightarrow$ intuición $\rightarrow$ formulación esencial $\rightarrow$ cómo funciona $\rightarrow$ aplicación con datos reales $\rightarrow$ evaluación $\rightarrow$ interpretación. Las derivaciones largas van al apéndice ``Para profundizar'' de cada sesión.
\item \textbf{Evaluación:} presentación de artículo 10\% \quad Problem Set 1 15\% \quad Problem Set 2 15\% \quad Proyecto final 60\%.
\end{enumerate}
\end{frame}

%%=================================================
\begin{frame}{El mapa de las 9 sesiones}
\begin{center}
\footnotesize
\begin{tabular}{cll}
\hline
\# & Tema & ISL \\
\hline
1 & Predecir, no explicar: el marco del aprendizaje estadístico \quad $\leftarrow$ \textbf{hoy} & 1--2 \\
2 & La regresión lineal como máquina de predecir; evaluar una predicción & 3, 7 \\
3 & Clasificación: de la logística a las métricas que deciden & 4 \\
4 & Remuestreo, validación y el \emph{pipeline} honesto & 5 \\
5 & Selección de variables y regularización & 6 \\
6 & Árboles, bosques y \emph{boosting} & 8 \\
7 & Abrir la caja negra: importancia, dependencia parcial y SHAP & --- \\
8 & Redes neuronales: qué son, cuándo usarlas y una aplicación & 10 \\
9 & Estructura sin $y$, la frontera causal $+$ presentaciones finales & 12 \\
\hline
\end{tabular}
\end{center}

\vspace{0.25cm}
\small
\begin{itemize}\itemsep3pt
  \item \textbf{Problem Set 1} cubre las sesiones 1--5; \textbf{Problem Set 2} cubre las sesiones 6--9.
  \item Desde la sesión 2, cada clase cierra con la \textbf{presentación de un artículo} por un estudiante (15 minutos: qué se predice, con qué datos y método, cómo se evalúa, cómo se interpreta).
  \item El \textbf{proyecto final} corre en paralelo durante todo el semestre.
\end{itemize}
\end{frame}

%%=================================================
\begin{frame}{Los datos que nos acompañan todo el semestre}
\small
Tres conjuntos de datos reales, reutilizados de sesión en sesión, para que el esfuerzo se vaya a los métodos y no a entender datos nuevos cada semana:
\pause
\begin{enumerate}\itemsep6pt
  \item \textbf{Mesas electorales de Colombia, 2006--2022.} Resultados por mesa (voto por el ganador, participación) enlazados con características censales del entorno (estrato, servicios, edades, composición étnica) y luces nocturnas. \emph{Regresión, clasificación, validación espacial, árboles e interpretación.} \pause
  \item \textbf{Precios de vivienda en Bogotá.} Anuncios con características físicas y de localización. \emph{Regresión, flexibilidad, regularización y árboles.} \pause
  \item \textbf{Imágenes satelitales y luces nocturnas.} \emph{La aplicación de redes neuronales.}
\end{enumerate}
\pause

\vspace{0.2cm}
\begin{block}{La pregunta que abre el curso}
¿Qué tan predecible es una elección a partir de las características \emph{estables} de los lugares? La respondemos por etapas: hoy con una recta, en la sesión 6 con un bosque, en la sesión 7 abriendo el bosque.
\end{block}
\end{frame}

%%=================================================
\begin{frame}{Nivelación en R y política de IA}
\small
\textbf{Nivelación en R} (obligatoria, autónoma, por fuera del horario de clase):
\begin{itemize}\itemsep4pt
  \item Introducción a R:
  \href{https://eduard-martinez.github.io/blog/intro-R/index.html}{Disponible \textcolor{blue}{[aquí]}}
  \item Manipulación de datos con \texttt{dplyr}:
  \href{https://eduard-martinez.github.io/blog/dplyr/index.html}{Disponible \textcolor{blue}{[aquí]}}
  \item Lectura y escritura de datos:
  \href{https://eduard-martinez.github.io/blog/read_and_write_data/index.html}{Disponible \textcolor{blue}{[aquí]}}
  \item Gestión y organización de datos:
  \href{https://eduard-martinez.github.io/blog/data_managment/index.html}{Disponible \textcolor{blue}{[aquí]}}
  \item Buenas prácticas de modelamiento con \texttt{tidymodels}: material en \textbf{Intu}.
  \item Para quienes opten por la vía Python, se ofrece material equivalente de referencia.
\end{itemize}
\pause

\vspace{0.2cm}
\begin{block}{Política de IA}
Se permite el uso de herramientas de IA en cualquier componente del curso, siempre que se \textbf{declare explícitamente} qué se usó y para qué.
\end{block}
\end{frame}


%%#################################################
%%#################################################
%% PARTE 2: DE BETA A Y
%%#################################################
%%#################################################
\section{De $\beta$ a $y$: dos preguntas para los mismos datos}
\begin{frame}[noframenumbering]
\tableofcontents[currentsection]
\end{frame}

%%=================================================
\begin{frame}{Los mismos datos, dos preguntas}
\framesubtitle{Lo que cambia respecto a Inferencia Causal}
\small
Tomen el panel de mesas electorales: voto por el ganador ($y$) y composición étnica del entorno ($x$).
\pause
\vspace{0.15cm}
\begin{columns}[T]
\begin{column}{0.48\textwidth}
\textbf{Pregunta causal (su otro curso)}
\begin{itemize}\itemsep3pt
  \item ¿Cuánto \emph{cambia} el voto si cambia la composición étnica?
  \item Objeto: un parámetro $\beta$.
  \item Enemigo: el sesgo (selección, omitidas).
  \item Se valida con un \textbf{diseño} de identificación.
\end{itemize}
\end{column}
\begin{column}{0.48\textwidth}
\textbf{Pregunta predictiva (este curso)}
\begin{itemize}\itemsep3pt
  \item ¿Qué tan bien podemos \emph{anticipar} el voto de una mesa que no hemos visto, y de qué depende?
  \item Objeto: la función completa $\hat{f}(x)$ y su error.
  \item Enemigo: el sobreajuste (varianza).
  \item Se valida con \textbf{datos nuevos}.
\end{itemize}
\end{column}
\end{columns}
\pause

\vspace{0.3cm}
\begin{block}{La frase que vamos a repetir todo el semestre}
Un coeficiente significativo no es una predicción, y la importancia de una variable en un modelo predictivo no es un efecto causal. Son preguntas distintas, con herramientas y criterios de éxito distintos.
\end{block}
\end{frame}

%%=================================================
\begin{frame}{¿Y qué es machine learning?}
\framesubtitle{Un cambio de paradigma: de estimar $\beta$ a predecir $y$}
\small
\begin{itemize}\itemsep5pt
  \item Una rama de la estadística y la computación que desarrolla \textbf{algoritmos para predecir} un resultado $y$ a partir de variables observables $x$. \pause
  \item La parte de ``aprendizaje'': \textbf{no especificamos} de antemano la forma exacta en que la máquina debe predecir $y$ a partir de $x$; eso queda como un problema empírico que el algoritmo resuelve con los datos. \pause
  \item El enfoque es pragmático: nos abstraemos del modelo subyacente.
  \begin{center}
  \emph{``Whatever works, works\ldots''}
  \end{center} \pause
  \item \textbf{¿En qué se diferencia de la econometría que ya conocen?}
  \begin{itemize}\itemsep3pt
    \item Econometría clásica: estimar \textbf{parámetros} ($\beta$) con propiedades deseables e interpretarlos.
    \item ML: producir \textbf{predicciones} ($\hat{y}$) que funcionen bien en datos \textbf{nuevos}, y luego entender de qué dependen.
  \end{itemize}
\end{itemize}
\end{frame}

%%=================================================
\begin{frame}{La primera gran victoria del ML: Google Flu Trends}
\small
\textbf{Contexto (2008):} en EE.UU., la vigilancia de la gripe dependía de los reportes que los centros médicos envían a la CDC.
\begin{itemize}\itemsep4pt
  \item La CDC agrega los reportes por región y a nivel nacional.
  \item Todo el proceso tomaba $\approx$ \textbf{10 días} $\rightarrow$ demasiado para una epidemia. \pause
\end{itemize}

\vspace{0.2cm}
\textbf{La propuesta de Google} (Ginsberg et al., 2009, \emph{Nature}):
\begin{itemize}\itemsep4pt
  \item Cruzar los datos oficiales con las \textbf{búsquedas} relacionadas con la gripe.
  \item Evaluaron $\approx$ \textbf{50 millones} de términos candidatos y $\approx$ \textbf{450 millones} de modelos; el modelo final usaba 45 términos. \pause
  \item Resultado: predicciones de intensidad de gripe casi en \textbf{tiempo real}, para cualquier región.
  \item A Google le tomaba \textbf{1 día} lo que a la CDC \textbf{10}.
\end{itemize} \pause

\vspace{0.2cm}
\textbf{Sin teoría epidemiológica, sin encuestas, sin modelo estructural.} ¿Convencidos?
\end{frame}

%%=================================================
\begin{frame}{\ldots y su primera gran derrota}
\framesubtitle{``The Parable of Google Flu'' (Lazer et al., 2014, Science)}
\small
\begin{itemize}\itemsep5pt
  \item En la temporada 2012--2013, Google Flu Trends predijo \textbf{más del doble} de los casos que efectivamente ocurrieron. \pause
  \item De hecho, sobreestimó la prevalencia en \textbf{100 de las 108 semanas} posteriores a agosto de 2011. \pause
  \item \textbf{¿Qué salió mal?}
  \begin{itemize}\itemsep3pt
    \item Con 50 millones de términos, muchos correlacionan con la gripe \textbf{por azar} (p.\,ej., búsquedas estacionales de baloncesto).
    \item El comportamiento de búsqueda \textbf{cambia con el tiempo} (y Google cambia su buscador): el modelo se entrenó en un mundo que dejó de existir.
    \item ``\emph{Big data hubris}'': la idea implícita de que el tamaño de los datos sustituye a la medición y a la validación cuidadosas.
  \end{itemize} \pause
\end{itemize}

\begin{block}{La lección que organiza este curso}
Predecir bien \textbf{dentro} de la muestra es fácil; lo difícil es \textbf{generalizar} a datos nuevos. Todo el aparato de hoy (riesgo, sobreajuste, sesgo--varianza) existe para atacar ese problema.
\end{block}
\end{frame}

%%=================================================
\begin{frame}{¿Paraguas o danza de la lluvia?}
\framesubtitle{Kleinberg, Ludwig, Mullainathan \& Obermeyer (2015): prediction policy problems}
\small
No toda pregunta de política necesita un efecto causal. Comparen: \pause
\begin{itemize}\itemsep5pt
  \item \textbf{¿Hago la danza de la lluvia?} Solo importa si la danza \emph{causa} lluvia $\rightarrow$ problema \textbf{causal} ($\beta$).
  \item \textbf{¿Llevo paraguas?} Solo importa si va a llover, no por qué llueve $\rightarrow$ problema de \textbf{predicción} ($\hat{y}$). \pause
\end{itemize}

\vspace{0.2cm}
\textbf{Ejemplos reales de problemas de predicción en política pública:}
\begin{itemize}\itemsep4pt
  \item \textbf{Cirugías fútiles:} miles de reemplazos de cadera/rodilla se hacen a pacientes que fallecen dentro del año siguiente; predecir mortalidad permitiría reasignar el gasto. \pause
  \item \textbf{Libertad provisional} (Kleinberg et al., 2018, QJE): jueces deciden con una predicción implícita de fuga o reincidencia. Un modelo de ML podría reducir el crimen $\approx 25\%$ sin más presos, o la población carcelaria $\approx 42\%$ sin más crimen. \pause
  \item \textbf{Focalización:} el Sisbén es, en esencia, un \textbf{clasificador} de hogares --- predice pobreza estructural a partir de características observables.
\end{itemize}
\end{frame}

%%=================================================
\begin{frame}{$\hat{y}$-problems vs.\ $\beta$-problems}
\framesubtitle{Mullainathan \& Spiess (2017)}
\small
Econometría I--II les enseñó a estimar $\beta$ creíbles. Este curso optimiza $\hat{y}$:

\vspace{0.2cm}
\begin{center}
\footnotesize
\begin{tabular}{lll}
\hline
 & \textbf{Inferencia causal} & \textbf{Predicción (ML)} \\
\hline
Objeto de interés & $\hat{\beta}$ & $\hat{y}$ \\
Pregunta típica & ¿qué pasa si intervengo? & ¿qué valor tomará $y$? \\
Enemigo principal & sesgo (selección, omitidas) & sobreajuste (varianza) \\
Validación & diseño $+$ supuestos & desempeño \emph{out-of-sample} \\
Ejemplo & ¿la beca mejora el Saber Pro? & ¿qué hogar caerá en pobreza? \\
\hline
\end{tabular}
\end{center}
\pause

\vspace{0.3cm}
\begin{itemize}\itemsep3pt
  \item Son preguntas \textbf{distintas} que exigen herramientas distintas: un modelo puede predecir muy bien con coeficientes sin interpretación causal, y viceversa.
  \item \textbf{No son rivales:} en la sesión 9 veremos cómo la buena predicción se pone al servicio de la inferencia causal.
\end{itemize}
\end{frame}

%%=================================================
\begin{frame}{Significativo no es predictivo}
\framesubtitle{Kim \& Zilinsky (2024): efectos marginales crecientes, capacidad predictiva estable}
\small
\begin{columns}[T]
\begin{column}{0.5\textwidth}
\begin{center}
\begin{tikzpicture}[scale=0.62]
  % panel izquierdo: efecto marginal
  \draw[->] (0,0) -- (4.4,0) node[below, font=\scriptsize] {año};
  \draw[->] (0,0) -- (0,3.2) node[above, font=\scriptsize] {efecto marginal};
  \draw[blue, very thick] plot[smooth] coordinates {(0.3,0.9) (1.3,1.1) (2.3,1.5) (3.3,2.2) (4.1,2.8)};
  \node[font=\scriptsize, blue] at (2.2,2.75) {educación $\uparrow$};
  \node[font=\scriptsize] at (2.2,-0.75) {(a) Coeficientes};
\end{tikzpicture}
\end{center}
\end{column}
\begin{column}{0.5\textwidth}
\begin{center}
\begin{tikzpicture}[scale=0.62]
  % panel derecho: accuracy plana
  \draw[->] (0,0) -- (4.4,0) node[below, font=\scriptsize] {año};
  \draw[->] (0,0) -- (0,3.2) node[above, font=\scriptsize] {accuracy};
  \fill[red!12] plot[smooth] coordinates {(0.3,1.95) (1.3,2.05) (2.3,1.9) (3.3,2.1) (4.1,2.0)} -- plot[smooth] coordinates {(4.1,1.6) (3.3,1.7) (2.3,1.5) (1.3,1.65) (0.3,1.55)} -- cycle;
  \draw[red, very thick] plot[smooth] coordinates {(0.3,1.75) (1.3,1.85) (2.3,1.7) (3.3,1.9) (4.1,1.8)};
  \draw[dashed, gray] (0,1.25) -- (4.3,1.25);
  \node[font=\scriptsize, gray] at (3.5,0.9) {azar (50\%)};
  \node[font=\scriptsize, red] at (2.2,2.55) {$\approx 64\%$ fuera de muestra};
  \node[font=\scriptsize] at (2.2,-0.75) {(b) Predicción};
\end{tikzpicture}
\end{center}
\end{column}
\end{columns}
\pause
\vspace{0.1cm}
\begin{itemize}\itemsep3pt
  \item Con solo demografía (edad, género, raza, educación, ingreso), \emph{random forests} y logit predicen el voto presidencial en EE.UU.\ con \textbf{63,9\%} de acierto fuera de muestra, estable de 1952 a 2020.
  \item Al mismo tiempo, los efectos marginales de educación crecen y son ``significativos''. Ambas cosas son ciertas: los coeficientes describen asociaciones; la predicción conjunta es otra cosa. \pause
  \item \textbf{Esquema} a partir de las figuras 2 y 3 del paper. Este artículo es el candidato natural para la primera presentación estudiantil.
\end{itemize}
\end{frame}

%%=================================================
\begin{frame}{La predictibilidad como cantidad de interés}
\framesubtitle{Gelvez, Cardiles, Martínez-González \& Muñoz (2026): ¿qué tan predecible es una elección?}
\small
La misma idea, con nuestros datos: ¿cuánto de una elección se puede anticipar a partir de las características estables de los lugares?
\pause
\begin{itemize}\itemsep4pt
  \item Modelos entrenados con el 80\% de las mesas y evaluados en el 20\% que \textbf{nunca vieron}, elección por elección.
  \item ¿Ganó el eventual presidente en la mesa? AUC entre \textbf{0,86 y 0,93}. ¿Qué proporción de votos obtuvo? $R^2$ fuera de muestra entre \textbf{0,48 y 0,68}. \pause
  \item La predictibilidad \textbf{no crece de forma monótona}: 2022 es la elección más ``legible'' desde el territorio aunque su resultado fue políticamente novedoso.
  \item Las mismas características (etnicidad, edades, servicios) siguen siendo útiles, pero \textbf{cambian de signo} según la coalición ganadora: el mapa sigue ahí; quien lo aprovecha cambia. \pause
\end{itemize}

\begin{block}{Qué nos enseña sobre el curso}
El desempeño fuera de muestra es el \emph{estadístico}; abrir el modelo (sesión 7) es lo que le da sentido económico. Ninguna de las dos cosas es un efecto causal.
\end{block}
\end{frame}


%%#################################################
%%#################################################
%% PARTE 3: EL MARCO DEL APRENDIZAJE ESTADÍSTICO
%%#################################################
%%#################################################
\section{El marco del aprendizaje estadístico}
\begin{frame}[noframenumbering]
\tableofcontents[currentsection]
\end{frame}

%%=================================================
\begin{frame}{El setup}
\small
\textbf{Datos:} observamos $n$ pares $(x_i, y_i)$, extraídos de forma i.i.d.\ de una distribución conjunta $P(X, Y)$ \textbf{desconocida}.
\begin{itemize}\itemsep3pt
  \item $Y$: variable de respuesta (salario, impago, consumo, precio, voto).
  \item $X = (X_1, \ldots, X_p)$: predictores o \emph{features}.
\end{itemize}
\pause

\vspace{0.2cm}
\textbf{Descomposición fundamental:} siempre podemos escribir
\[ Y \;=\; f(X) \;+\; \varepsilon, \qquad \text{con } f(x) \equiv \E[Y \mid X = x] \]
\begin{itemize}\itemsep3pt
  \item $f$: la \textbf{función de regresión} --- la parte sistemática de $Y$.
  \item $\varepsilon \equiv Y - f(X)$: lo que $X$ no explica. \textbf{Por construcción} $\E[\varepsilon \mid X] = 0$, con $\Var(\varepsilon) = \sigma^2$.
\end{itemize}
\pause

\vspace{0.2cm}
\begin{block}{La meta del aprendizaje supervisado}
Usar los datos para construir $\hat{f} \approx f$. En econometría nos importaba un \textbf{parámetro dentro} de $f$; aquí nos importa la \textbf{función completa}, porque con ella predecimos.
\end{block}
\end{frame}

%%=================================================
\begin{frame}{¿Qué significa exactamente ``aprender''?}
\small
\textbf{Aprender} $=$ elegir $\hat{f}$ dentro de una \textbf{clase de funciones} $\mathcal{F}$ usando los datos:
\begin{itemize}\itemsep3pt
  \item $\mathcal{F}$ puede ser: rectas, polinomios, árboles, redes neuronales\ldots
  \item Cada sesión de este curso es, en el fondo, una clase $\mathcal{F}$ distinta.
\end{itemize}
\pause

\vspace{0.2cm}
\textbf{¿Y qué es una ``buena'' $\hat{f}$?} Necesitamos una \textbf{función de pérdida} $L(y, \hat{f}(x))$ que mida el costo de equivocarse:
\begin{itemize}\itemsep4pt
  \item \textbf{Cuadrática:} $L = (y - \hat{f}(x))^2$. Castiga fuerte los errores grandes; estándar en regresión.
  \item \textbf{Absoluta:} $L = |y - \hat{f}(x)|$. Más robusta a valores extremos.
  \item \textbf{0--1:} $L = \1\{y \neq \hat{f}(x)\}$. Para clasificación (sesión 3).
\end{itemize}
\pause

\vspace{0.2cm}
\begin{block}{La elección de $L$ es una decisión económica}
¿Qué error cuesta más? No es lo mismo subestimar la demanda (ventas perdidas) que sobreestimarla (inventario muerto); ni un falso positivo que un falso negativo al focalizar un subsidio.
\end{block}
\end{frame}

%%=================================================
\begin{frame}{El riesgo de generalización}
\small
\textbf{Definición.} El \textbf{riesgo} (o error de generalización) de un predictor $f$ es su pérdida esperada sobre un dato \textbf{nuevo} extraído de $P$:
\[ R(f) \;=\; \E\big[ L(Y, f(X)) \big] \]
\pause

\vspace{0.1cm}
\textbf{¿Cuál es el mejor predictor posible bajo pérdida cuadrática?} La \textbf{media condicional}:
\[ f^\ast(x) \;=\; \E[Y \mid X = x] \;=\; \argmin_{f} R(f) \]
\begin{itemize}\itemsep3pt
  \item Intuición: en cada $x$, la constante que minimiza el error cuadrático esperado es el promedio de $Y$ en ese $x$; cualquier otra predicción añade el cuadrado de su distancia a ese promedio (la cuenta, de tres líneas, está en el apéndice).
\end{itemize}
\pause
\vspace{0.2cm}
\begin{block}{Por qué la regresión es el punto de partida}
La regresión aproxima la media condicional. Por eso es el primer método del curso (sesión 2) y el \emph{baseline} contra el que se mide todo lo demás.
\end{block}
\end{frame}

%%=================================================
\begin{frame}{Error reducible y error irreducible}
\small
\textbf{Lo que ni el mejor predictor evita:} para cualquier $\hat{f}$ y cualquier $x$, con $Y = f(x) + \varepsilon$,
\[ \E\big[(Y - \hat{f}(x))^2\big] \;=\; \underbrace{\big(f(x) - \hat{f}(x)\big)^2}_{\text{\textbf{reducible}}} \;+\; \underbrace{\sigma^2}_{\text{\textbf{irreducible}}} \]
\pause
\begin{itemize}\itemsep5pt
  \item \textbf{Reducible:} qué tan lejos está nuestra $\hat{f}$ de la verdadera $f$. Es lo que un mejor método, más datos o mejores variables pueden mejorar.
  \item \textbf{Irreducible} ($\sigma^2$): ni el modelo perfecto ($\hat{f} = f$) lo elimina. Es el techo de lo aprendible con \emph{estos} predictores. (La cuenta, de dos líneas, en el apéndice.)
\end{itemize}
\pause
\vspace{0.2cm}
\begin{block}{Mensaje}
Si el desempeño de un modelo parece ``malo'', la causa puede ser $\sigma^2$ alto: no siempre hay un mejor algoritmo; a veces faltan mejores \textbf{datos}. Kim \& Zilinsky lo ilustran: con solo demografía, el techo de la predictibilidad del voto es bajo, y ningún algoritmo lo sube.
\end{block}
\end{frame}

%%=================================================
\begin{frame}{Paramétrico o no, flexible o interpretable}
\framesubtitle{Dos decisiones que toma todo método (ISL, secciones 2.1.2--2.1.3)}
\small
\begin{columns}[T]
\begin{column}{0.44\textwidth}
\textbf{Cómo estimar $f$:}
\begin{itemize}\itemsep2pt
  \item \textbf{Paramétrico:} fijar la forma ($f = \beta_0 + \beta_1 x + \cdots$) y estimar pocos números. Estable; puede estar \emph{mal especificado}.
  \item \textbf{No paramétrico:} dejar que los datos dibujen la forma ($k$-NN, \emph{splines}, árboles). Flexible; necesita \emph{muchos} datos.
\end{itemize}
\end{column}
\begin{column}{0.56\textwidth}
\begin{center}
\begin{tikzpicture}[scale=0.5]
  \draw[->] (0,0) -- (6.2,0) node[below, font=\scriptsize] {Flexibilidad $\rightarrow$};
  \draw[->] (0,0) -- (0,4.8) node[above, font=\scriptsize] {Interpretabilidad $\uparrow$};
  \node[font=\scriptsize, fill=blue!10, rounded corners, inner sep=2pt] at (0.9,4.2) {Lasso};
  \node[font=\scriptsize, fill=blue!10, rounded corners, inner sep=2pt] at (1.9,3.5) {OLS};
  \node[font=\scriptsize, fill=blue!10, rounded corners, inner sep=2pt] at (2.7,2.75) {\emph{splines}, GAM};
  \node[font=\scriptsize, fill=blue!10, rounded corners, inner sep=2pt] at (3.4,2.0) {árboles};
  \node[font=\scriptsize, fill=blue!10, rounded corners, inner sep=2pt] at (4.6,1.25) {RF, \emph{boosting}};
  \node[font=\scriptsize, fill=blue!10, rounded corners, inner sep=2pt] at (5.4,0.5) {redes};
  \node[font=\scriptsize, gray] at (3.1,-0.85) {Esquema tipo ISL, figura 2.7};
\end{tikzpicture}
\end{center}
\end{column}
\end{columns}
\pause
\vspace{0.1cm}
\begin{block}{El precio de la flexibilidad, y por qué existe la sesión 7}
Bajar por la diagonal compra precisión y pierde lectura (y arriesga ajustar el ruido). Importancia de variables, dependencia parcial y SHAP recuperan parte de la lectura.
\end{block}
\end{frame}

%%=================================================
\begin{frame}{El riesgo empírico y su trampa}
\small
\textbf{Problema:} $R(f) = \E[L(Y, f(X))]$ depende de $P$, que no conocemos.
\pause

\vspace{0.1cm}
\textbf{Solución natural:} estimarlo con el promedio muestral --- el \textbf{riesgo empírico}:
\[ \hat{R}(f) \;=\; \frac{1}{n} \sum_{i=1}^{n} L\big(y_i, f(x_i)\big) \]
y elegir el modelo minimizándolo (\emph{empirical risk minimization}):
\[ \hat{f} \;=\; \argmin_{f \in \mathcal{F}} \; \hat{R}(f) \]
\pause

\textbf{La trampa:} evaluar $\hat{f}$ con los \textbf{mismos datos} con que se eligió.
\begin{itemize}\itemsep4pt
  \item $\hat{R}(\hat{f})$ es \textbf{optimista}: $\hat{f}$ fue construida \emph{precisamente} para minimizarlo. Es presentar el examen con las respuestas en la mano. \pause
  \item Con $\mathcal{F}$ suficientemente flexible, $\hat{R}(\hat{f}) = 0$: basta \textbf{memorizar} los datos. Y un modelo que memoriza no ha aprendido nada generalizable.
\end{itemize}
\pause

\begin{block}{Regla de oro \#1}
El error de entrenamiento \textbf{no} estima el error de generalización.
\end{block}
\end{frame}

%%=================================================
\begin{frame}{Entrenamiento vs.\ prueba}
\small
La solución más simple: separar los datos \textbf{antes} de tocar cualquier modelo.

\vspace{0.3cm}
\begin{center}
\begin{tikzpicture}
  \draw[fill=blue!12] (0,0) rectangle (6.5,0.9);
  \node at (3.25,0.45) {\small Entrenamiento ($\approx 70$--$80\%$)};
  \draw[fill=red!12] (6.5,0) rectangle (9.2,0.9);
  \node at (7.85,0.45) {\small Prueba};
  \draw[-{Stealth}, thick] (3.25,-0.15) -- (3.25,-0.7);
  \node[font=\scriptsize] at (3.25,-0.95) {elegir y ajustar $\hat{f}$};
  \draw[-{Stealth}, thick] (7.85,-0.15) -- (7.85,-0.7);
  \node[font=\scriptsize] at (7.85,-0.95) {estimar $R(\hat{f})$, \textbf{una sola vez}};
\end{tikzpicture}
\end{center}
\pause

\vspace{0.3cm}
\begin{itemize}\itemsep4pt
  \item El conjunto de \textbf{prueba} simula los ``datos nuevos'': su error sí estima $R(\hat{f})$ honestamente.
  \item \textbf{Regla de oro \#2:} el conjunto de prueba se usa \textbf{al final y una sola vez}. Si lo usamos muchas veces para escoger el modelo, se convierte --- en la práctica --- en un segundo conjunto de entrenamiento. \pause
  \item Con esto ya podemos trabajar en las sesiones 2 y 3. ¿Y si la muestra es pequeña y partirla ``duele''? De eso se encarga la \textbf{validación cruzada} (sesión 4).
\end{itemize}
\end{frame}


%%#################################################
%%#################################################
%% PARTE 4: SOBREAJUSTE Y COMPLEJIDAD
%%#################################################
%%#################################################
\section{Sobreajuste y complejidad}
\begin{frame}[noframenumbering]
\tableofcontents[currentsection]
\end{frame}

%%=================================================
\begin{frame}{Un mismo conjunto de datos, tres modelos}
\small
Ajustemos polinomios de grado creciente a los mismos 9 puntos:

\vspace{0.2cm}
\begin{center}
\begin{tikzpicture}[scale=0.60]
% Panel A: grado 1
\begin{scope}[shift={(0,0)}]
  \draw[->] (0,0) -- (4.7,0);
  \draw[->] (0,0) -- (0,4.2);
  \foreach \px/\py in {0.2/1.0, 0.7/2.1, 1.2/2.5, 1.7/2.1, 2.2/2.0, 2.7/1.4, 3.2/1.9, 3.7/2.4, 4.2/3.4} \fill (\px,\py) circle (2.6pt);
  \draw[blue, very thick] (0.1,1.55) -- (4.4,2.70);
  \node[font=\footnotesize] at (2.3,-0.6) {Grado 1: \textbf{subajuste}};
\end{scope}
% Panel B: grado 3
\begin{scope}[shift={(5.9,0)}]
  \draw[->] (0,0) -- (4.7,0);
  \draw[->] (0,0) -- (0,4.2);
  \foreach \px/\py in {0.2/1.0, 0.7/2.1, 1.2/2.5, 1.7/2.1, 2.2/2.0, 2.7/1.4, 3.2/1.9, 3.7/2.4, 4.2/3.4} \fill (\px,\py) circle (2.6pt);
  \draw[blue, very thick] plot[smooth, tension=0.7] coordinates {(0.15,1.2) (1.2,2.45) (2.2,1.95) (2.7,1.55) (3.2,1.8) (4.2,3.35)};
  \node[font=\footnotesize] at (2.3,-0.6) {Grado 3: $\approx f$};
\end{scope}
% Panel C: grado 15
\begin{scope}[shift={(11.8,0)}]
  \draw[->] (0,0) -- (4.7,0);
  \draw[->] (0,0) -- (0,4.2);
  \foreach \px/\py in {0.2/1.0, 0.7/2.1, 1.2/2.5, 1.7/2.1, 2.2/2.0, 2.7/1.4, 3.2/1.9, 3.7/2.4, 4.2/3.4} \fill (\px,\py) circle (2.6pt);
  \draw[blue, very thick] plot[smooth, tension=0.85] coordinates {(0.2,1.0) (0.45,2.7) (0.7,2.1) (0.95,1.6) (1.2,2.5) (1.45,2.95) (1.7,2.1) (1.95,1.45) (2.2,2.0) (2.45,2.6) (2.7,1.4) (2.95,0.8) (3.2,1.9) (3.45,2.9) (3.7,2.4) (3.95,1.5) (4.2,3.4)};
  \node[font=\footnotesize] at (2.3,-0.6) {Grado 15: \textbf{sobreajuste}};
\end{scope}
\end{tikzpicture}
\end{center}
\pause

\vspace{0.2cm}
\begin{itemize}\itemsep3pt
  \item \textbf{Grado 1:} demasiado rígido --- no captura la forma de $f$ (error por \textbf{sesgo}).
  \item \textbf{Grado 15:} pasa por \emph{todos} los puntos --- ajusta el \textbf{ruido} $\varepsilon$, no la señal (error por \textbf{varianza}).
  \item ¿Cuál predice mejor un punto \textbf{nuevo}? El del medio. Ni el más simple ni el que ``mejor ajusta''.
\end{itemize}
\end{frame}

%%=================================================
\begin{frame}{La curva que resume el curso}
\small
\begin{center}
\begin{tikzpicture}[scale=0.80]
  \draw[->] (0,0) -- (7.8,0) node[below, font=\scriptsize] {Complejidad $\rightarrow$};
  \draw[->] (0,0) -- (0,4.5) node[above, font=\scriptsize] {Error};
  \draw[blue, very thick] plot[smooth] coordinates {(0.5,3.5) (1.4,2.2) (2.4,1.45) (3.4,1.0) (4.4,0.72) (5.4,0.52) (6.9,0.38)};
  \draw[red, very thick] plot[smooth] coordinates {(0.5,3.7) (1.4,2.5) (2.4,1.95) (3.2,1.72) (3.9,1.78) (5.0,2.3) (6.9,3.5)};
  \draw[dashed, gray, thick] (3.4,0) -- (3.4,4.0);
  \node[font=\scriptsize, gray] at (3.4,4.2) {óptimo};
  \node[font=\scriptsize, blue] at (6.0,0.9) {error de entrenamiento};
  \node[font=\scriptsize, red] at (5.5,3.45) {error de prueba};
  \node[font=\scriptsize] at (0.9,4.2) {subajuste};
  \node[font=\scriptsize] at (6.2,4.2) {sobreajuste};
\end{tikzpicture}
\end{center}
\pause

\begin{itemize}\itemsep3pt
  \item El error de \textbf{entrenamiento} cae \emph{monótonamente} con la complejidad: más flexibilidad siempre ajusta mejor la muestra.
  \item El error de \textbf{prueba} tiene forma de \textbf{U}: mejora hasta cierto punto y luego el sobreajuste lo destruye.
  \item \textbf{Todo el curso vive en esta figura:} cada método trae una ``perilla'' de complejidad y una forma de encontrar el mínimo de la curva roja.
\end{itemize}
\end{frame}

%%=================================================
\begin{frame}{¿Qué es exactamente la ``complejidad''?}
\small
Informalmente: la \textbf{capacidad} de la clase $\mathcal{F}$ para ajustar patrones arbitrarios.
\pause
\vspace{0.2cm}
\begin{itemize}\itemsep5pt
  \item En regresión polinómica: el \textbf{grado} del polinomio.
  \item En regresión múltiple: el \textbf{número de predictores} $p$.
  \item Más adelante en el curso, cada método tendrá su propia perilla:
  \begin{itemize}\itemsep3pt
    \item $k$-NN: el número de vecinos $k$ (sesiones 2--3).
    \item Regularización: el parámetro de penalización $\lambda$ (sesión 5).
    \item Árboles: la profundidad, el tamaño de las hojas y el número de árboles (sesión 6).
    \item Redes neuronales: capas, neuronas, \emph{dropout} (sesión 8).
  \end{itemize} \pause
\end{itemize}

\begin{block}{Idea clave}
Elegir el modelo es elegir \textbf{dónde pararse} sobre la curva U. Necesitamos: (i) una perilla de complejidad y (ii) un estimador honesto del error de prueba para girarla. Hoy tenemos la partición entrenamiento/prueba; la sesión 4 construye el estimador definitivo.
\end{block}
\end{frame}


%%#################################################
%%#################################################
%% PARTE 5: SESGO Y VARIANZA, CON INTUICIÓN
%%#################################################
%%#################################################
\section{Sesgo y varianza: por qué el error de prueba tiene forma de U}
\begin{frame}[noframenumbering]
\tableofcontents[currentsection]
\end{frame}

%%=================================================
\begin{frame}{El experimento mental}
\small
Queremos entender \emph{por qué} el error de prueba tiene forma de U. Fijemos un punto $x_0$ y pensemos qué pasaría si pudiéramos repetir el muestreo:
\pause
\begin{itemize}\itemsep5pt
  \item La muestra de entrenamiento $\mathcal{D}$ es \textbf{aleatoria}: con otros $n$ datos, obtendríamos \textbf{otra} $\hat{f}$.
  \item Por lo tanto, $\hat{f}(x_0)$ es una \textbf{variable aleatoria}: cambia de muestra a muestra. \pause
  \item Dos objetos naturales:
  \begin{itemize}\itemsep3pt
    \item $\E[\hat{f}(x_0)]$: la predicción \emph{promedio} entre muestras posibles.
    \item $\Var(\hat{f}(x_0))$: cuánto \emph{cambia} la predicción de muestra a muestra.
  \end{itemize} \pause
  \item Definimos el \textbf{sesgo} del método en $x_0$:
  \[ \mathrm{Sesgo}(x_0) \;=\; \E[\hat{f}(x_0)] - f(x_0) \]
\end{itemize}
\pause
\begin{block}{Pregunta}
Llega un dato nuevo $(x_0, Y_0)$ con $Y_0 = f(x_0) + \varepsilon_0$. ¿Cuál es el error cuadrático esperado $\E[(Y_0 - \hat{f}(x_0))^2]$, promediando sobre muestras $\mathcal{D}$ y sobre $\varepsilon_0$?
\end{block}
\end{frame}

%%=================================================
\begin{frame}{La descomposición sesgo--varianza}
\framesubtitle{Tres piezas, y solo dos dependen de nosotros}
\small
\textbf{Resultado} (la derivación, en dos pasos, está en el apéndice):
\[ \MSE(x_0) \;=\; \underbrace{\big(f(x_0) - \E[\hat{f}(x_0)]\big)^2}_{\text{\textbf{Sesgo}}^2} \;+\; \underbrace{\Var\big(\hat{f}(x_0)\big)}_{\text{\textbf{Varianza}}} \;+\; \underbrace{\sigma^2}_{\text{\textbf{Irreducible}}} \]
\pause
\vspace{0.1cm}
\begin{itemize}\itemsep5pt
  \item \textbf{Sesgo:} error sistemático del método. Un modelo rígido (la recta de grado 1) se equivoca siempre en la misma dirección, con cualquier muestra.
  \item \textbf{Varianza:} sensibilidad a la muestra particular. Un modelo flexible (grado 15) cambia por completo si cambian unos pocos puntos.
  \item $\sigma^2$: ruido que nadie puede eliminar; el techo de lo aprendible con estos predictores. \pause
\end{itemize}

\begin{block}{Por qué importa}
La complejidad \textbf{baja el sesgo y sube la varianza}. Como el error total es la suma, el óptimo es interior: esa es la curva U, y toda perilla del curso (grado, $k$, $\lambda$, profundidad, épocas) se mueve sobre ella.
\end{block}
\end{frame}

%%=================================================
\begin{frame}{Intuición: las cuatro dianas}
\small
Cada punto es la predicción $\hat{f}(x_0)$ obtenida con una muestra de entrenamiento distinta; el centro es el valor verdadero $f(x_0)$:

\vspace{0.2cm}
\begin{center}
\begin{tikzpicture}[scale=0.62]
% headers
\node[font=\footnotesize] at (0,1.75)  {\textbf{Varianza baja}};
\node[font=\footnotesize] at (4.4,1.75) {\textbf{Varianza alta}};
\node[font=\footnotesize, rotate=90] at (-1.95,0)  {\textbf{Sesgo bajo}};
\node[font=\footnotesize, rotate=90] at (-1.95,-4.4) {\textbf{Sesgo alto}};
% top-left: low bias low var
\begin{scope}[shift={(0,0)}]
  \draw[gray!70] (0,0) circle (1.3); \draw[gray!70] (0,0) circle (0.85); \fill[red!25] (0,0) circle (0.4); \draw[gray!70] (0,0) circle (0.4);
  \foreach \px/\py in {0.14/0.12, -0.18/0.06, 0.02/-0.19, -0.08/0.17, 0.19/-0.07} \fill[blue] (\px,\py) circle (3pt);
\end{scope}
% top-right: low bias high var
\begin{scope}[shift={(4.4,0)}]
  \draw[gray!70] (0,0) circle (1.3); \draw[gray!70] (0,0) circle (0.85); \fill[red!25] (0,0) circle (0.4); \draw[gray!70] (0,0) circle (0.4);
  \foreach \px/\py in {0.7/0.55, -0.75/0.4, 0.2/-0.9, -0.45/-0.65, 0.9/-0.15} \fill[blue] (\px,\py) circle (3pt);
\end{scope}
% bottom-left: high bias low var
\begin{scope}[shift={(0,-4.4)}]
  \draw[gray!70] (0,0) circle (1.3); \draw[gray!70] (0,0) circle (0.85); \fill[red!25] (0,0) circle (0.4); \draw[gray!70] (0,0) circle (0.4);
  \foreach \px/\py in {0.62/0.55, 0.78/0.68, 0.58/0.75, 0.72/0.48, 0.85/0.62} \fill[blue] (\px,\py) circle (3pt);
\end{scope}
% bottom-right: high bias high var
\begin{scope}[shift={(4.4,-4.4)}]
  \draw[gray!70] (0,0) circle (1.3); \draw[gray!70] (0,0) circle (0.85); \fill[red!25] (0,0) circle (0.4); \draw[gray!70] (0,0) circle (0.4);
  \foreach \px/\py in {0.95/0.25, 0.25/0.95, 1.1/0.85, 0.4/0.2, 0.75/1.15} \fill[blue] (\px,\py) circle (3pt);
\end{scope}
\end{tikzpicture}
\end{center}
\pause

\begin{itemize}\itemsep2pt
  \item \textbf{Sesgo:} ¿el \emph{centro} de la nube está sobre el blanco?
  \item \textbf{Varianza:} ¿qué tan \emph{dispersa} es la nube?
\end{itemize}
\end{frame}

%%=================================================
\begin{frame}{El trade-off sesgo--varianza}
\small
\begin{center}
\begin{tikzpicture}[scale=0.78]
  \draw[->] (0,0) -- (7.8,0) node[below, font=\scriptsize] {Complejidad $\rightarrow$};
  \draw[->] (0,0) -- (0,4.5) node[above, font=\scriptsize] {Error};
  \draw[blue, very thick] plot[smooth] coordinates {(0.5,3.2) (1.6,2.0) (2.6,1.2) (3.6,0.7) (4.6,0.42) (5.6,0.28) (6.8,0.2)};
  \draw[red, very thick] plot[smooth] coordinates {(0.5,0.15) (1.6,0.35) (2.6,0.65) (3.6,1.05) (4.6,1.6) (5.6,2.35) (6.8,3.3)};
  \draw[dashed, gray, thick] (0,0.5) -- (7.4,0.5);
  \node[font=\scriptsize, gray] at (6.7,0.72) {$\sigma^2$};
  \draw[black, ultra thick] plot[smooth] coordinates {(0.5,3.85) (1.6,2.85) (2.6,2.35) (3.6,2.25) (4.6,2.52) (5.6,3.13) (6.8,4.0)};
  \draw[dashed, gray] (3.7,0) -- (3.7,4.2);
  \node[font=\scriptsize, blue] at (1.5,1.5) {Sesgo$^2$};
  \node[font=\scriptsize, red] at (5.85,1.8) {Varianza};
  \node[font=\scriptsize] at (2.4,3.25) {$\MSE$ total};
\end{tikzpicture}
\end{center}
\pause

\begin{itemize}\itemsep3pt
  \item Más complejidad: el sesgo$^2$ \textbf{cae} (menos rigidez) pero la varianza \textbf{sube} (más sensibilidad a la muestra).
  \item El $\MSE$ total es la \textbf{suma} de las tres piezas: su mínimo es un \textbf{punto interior} --- ni el modelo más simple ni el más flexible.
  \item Esta es la explicación formal de la curva U de la sección anterior.
\end{itemize}
\end{frame}

%%=================================================
\begin{frame}{¿Por qué un economista aceptaría un estimador sesgado?}
\small
\textbf{El reflejo de Econometría I:} bajo Gauss--Markov, OLS es \textbf{MELI} --- el de mínima varianza\ldots\ \emph{entre los insesgados}. La insesgadez se impone como restricción \textbf{a priori}.
\pause

\vspace{0.2cm}
\textbf{La lógica de la predicción es distinta:} minimizamos el $\MSE$ total,
\[ \MSE \;=\; \text{Sesgo}^2 + \text{Varianza} + \sigma^2, \]
y a nadie le pagan por insesgadez si la varianza explota.
\pause

\vspace{0.2cm}
\textbf{¿Cuándo explota la varianza de OLS?} Con muchos predictores respecto a $n$, con predictores muy correlacionados, con muestras pequeñas o ruidosas. En esos casos, un estimador \textbf{levemente sesgado} pero mucho \textbf{más estable} predice mejor.
\pause

\vspace{0.2cm}
\begin{block}{Teaser de la sesión 5}
Ridge y Lasso hacen exactamente ese intercambio, con una perilla $\lambda$ que controla cuánto sesgo compramos a cambio de cuánta varianza. (De hecho, ni siquiera en estimación pura la insesgadez es sagrada: James \& Stein, 1961.)
\end{block}
\end{frame}


%%#################################################
%%#################################################
%% PARTE 6: EL MAPA DE TAREAS Y EL FLUJO DE TRABAJO
%%#################################################
%%#################################################
\section{El mapa de tareas y el flujo de trabajo del curso}
\begin{frame}[noframenumbering]
\tableofcontents[currentsection]
\end{frame}

%%=================================================
\begin{frame}{Supervisado y no supervisado: el mapa de tareas}
\small
\begin{center}
\footnotesize
\setlength{\tabcolsep}{4pt}
\begin{tabular}{lllll}
\hline
Tarea & ¿Hay $y$? & Pregunta & Ejemplo & Sesión \\
\hline
Regresión & Sí (continua) & ¿cuánto? & precio de vivienda & 2, 4--8 \\
Clasificación & Sí (discreta) & ¿qué clase? & ¿ganó en la mesa? & 3, 6--8 \\
\emph{Clustering} & No & ¿qué grupos? & segmentar hogares & 9 \\
Reducción de dim. & No & ¿qué estructura? & índice socioeconómico & 5, 9 \\
\hline
\end{tabular}
\end{center}
\pause

\vspace{0.3cm}
\begin{itemize}\itemsep4pt
  \item \textbf{Supervisado:} conocemos la respuesta $y$ en los datos de entrenamiento; el algoritmo aprende el mapa $x \rightarrow y$. La ``supervisión'' es la $y$ observada.
  \item \textbf{No supervisado:} no hay $y$; buscamos estructura en los $x$ (grupos, factores latentes, dimensiones). Sin $y$ tampoco hay error de prueba: validar es otra cosa (sesión 9).
  \item Las imágenes (sesión 8) combinan ambos mundos: primero se \emph{representan}, luego se predice con ellas.
\end{itemize}
\end{frame}

%%=================================================
\begin{frame}{Las dos culturas del modelamiento (Breiman, 2001)}
\small
Toda pregunta empírica parte del mismo diagrama: la naturaleza asocia insumos con resultados.

\vspace{0.2cm}
\begin{center}
\begin{tikzpicture}[node distance=1.1cm, every node/.style={font=\small}]
  \node (x) [draw, rounded corners, minimum height=0.8cm] {$x$};
  \node (nat) [draw, fill=gray!20, rounded corners, right=1.1cm of x, minimum height=0.8cm] {naturaleza};
  \node (y) [draw, rounded corners, right=1.1cm of nat, minimum height=0.8cm] {$y$};
  \draw[-{Stealth}, thick] (x) -- (nat);
  \draw[-{Stealth}, thick] (nat) -- (y);
\end{tikzpicture}
\end{center}
\pause

\begin{itemize}\itemsep5pt
  \item \textbf{Cultura 1 --- modelamiento de datos:} asumir un modelo estocástico para la caja (lineal, logit\ldots), estimar sus parámetros y validar con supuestos y bondad de ajuste. En 2001, el $\approx 98\%$ de los estadísticos (y casi toda la econometría). \pause
  \item \textbf{Cultura 2 --- modelamiento algorítmico:} tratar la caja como \textbf{desconocida}; buscar cualquier $f(x)$ que prediga bien y validar con error \emph{out-of-sample}. Entonces el $\approx 2\%$; hoy, la industria entera. \pause
\end{itemize}

\begin{block}{¿Dónde se para este curso?}
En la frontera, que ya se disolvió: la econometría adoptó los algoritmos y el ML aprendió causalidad (sesión 9). Pero adoptamos la \textbf{disciplina de validación} de la cultura 2: la prueba de un modelo es predecir datos que no vio.
\end{block}
\end{frame}

%%=================================================
\begin{frame}{Aplicación guiada: una elección y una recta}
\framesubtitle{El primer ejercicio del curso, sobre el panel de mesas electorales}
\small
\begin{enumerate}\itemsep2pt
  \item \textbf{Pregunta:} ¿cuánto del \emph{vote share} del ganador en 2022 se anticipa con las características censales del entorno de cada mesa? \pause
  \item \textbf{Partición:} 80\% de las mesas para entrenar y 20\% para probar, estratificando por deciles de $y$. El 20\% queda \emph{bajo llave}. \pause
  \item \textbf{Modelo:} una regresión lineal con los cuatro bloques de predictores (estrato, servicios y densidad, edades, composición étnica). \pause
  \item \textbf{Evaluación:} $R^2$ en entrenamiento y en prueba; luego añadimos interacciones y polinomios y vemos cómo se mueven \emph{los dos}. \pause
  \item \textbf{Lectura:} ¿qué mesas predecimos peor: las de apoyo muy alto, muy bajo, las rurales? Esa pregunta abre la sesión 2.
\end{enumerate}
\pause

\begin{block}{La vara con la que nos mediremos}
Con \emph{random forests}, este problema alcanza un $R^2$ fuera de muestra de \textbf{0,66} en la primera vuelta de 2022 (Gelvez et al., 2026). Cuánto recorre una recta y cuánto exige flexibilidad es la historia de las sesiones 2 a 6.
\end{block}
\end{frame}

%%=================================================
\begin{frame}{El flujo de trabajo del curso en una diapositiva}
\small
\begin{center}
\begin{tikzpicture}[node distance=0.32cm, every node/.style={font=\scriptsize, align=center}]
  \node (d) [draw, rounded corners, fill=gray!12, minimum height=0.9cm, text width=1.2cm] {datos};
  \node (p) [draw, rounded corners, fill=blue!10, minimum height=0.9cm, text width=1.35cm, right=of d] {partición\\ train / test};
  \node (r) [draw, rounded corners, fill=blue!10, minimum height=0.9cm, text width=1.5cm, right=of p] {preprocesar\\ \emph{dentro} del pipeline};
  \node (m) [draw, rounded corners, fill=blue!10, minimum height=0.9cm, text width=1.35cm, right=of r] {modelo $+$\\ validación};
  \node (e) [draw, rounded corners, fill=red!10, minimum height=0.9cm, text width=1.4cm, right=of m] {evaluar en\\ test (una vez)};
  \node (i) [draw, rounded corners, fill=red!10, minimum height=0.9cm, text width=1.35cm, right=of e] {interpretar\\ y reportar};
  \draw[-{Stealth}, thick] (d) -- (p);
  \draw[-{Stealth}, thick] (p) -- (r);
  \draw[-{Stealth}, thick] (r) -- (m);
  \draw[-{Stealth}, thick] (m) -- (e);
  \draw[-{Stealth}, thick] (e) -- (i);
  \node[font=\scriptsize, gray, below=0.15cm of p] {S1};
  \node[font=\scriptsize, gray, below=0.15cm of r] {S4};
  \node[font=\scriptsize, gray, below=0.15cm of m] {S2--S8};
  \node[font=\scriptsize, gray, below=0.15cm of e] {S2--S3};
  \node[font=\scriptsize, gray, below=0.15cm of i] {S7};
\end{tikzpicture}
\end{center}
\pause

\vspace{0.15cm}
Cada sesión cambia una pieza y deja las demás fijas:
\begin{itemize}\itemsep3pt
  \item \textbf{S2--S3:} qué modelo de referencia usar y cómo medir el error (regresión y clasificación). \pause
  \item \textbf{S4:} cómo validar sin engañarse (CV, \emph{leakage}, estructura espacial y temporal). \pause
  \item \textbf{S5--S6, S8:} clases $\mathcal{F}$ cada vez más flexibles, cada una con su perilla. \pause
  \item \textbf{S7:} cómo abrir la caja; \textbf{S9:} qué hacer sin $y$ y cuándo la predicción sirve para estimar $\beta$.
\end{itemize}
\pause
\begin{block}{Mensaje}
Si entienden la sesión de hoy --- riesgo, sobreajuste, sesgo--varianza, entrenamiento y prueba --- entienden \emph{por qué existe} cada pieza del resto del curso.
\end{block}
\end{frame}

%%=================================================
\begin{frame}{Recap}
\small
\begin{enumerate}\itemsep7pt
  \item \textbf{$Y$, no $\beta$:} predecir y explicar son preguntas distintas. Un coeficiente significativo no es una predicción (Kim \& Zilinsky), y muchos problemas de política son de predicción.
  \item El objetivo del ML supervisado es minimizar el \textbf{riesgo de generalización} $R(f) = \E[L(Y, f(X))]$, no el error en la muestra; el mejor predictor cuadrático es la media condicional.
  \item \textbf{Regla de oro \#1:} el error de entrenamiento es optimista. \textbf{Regla de oro \#2:} el conjunto de prueba se usa al final y una sola vez.
  \item \textbf{Descomposición fundamental:} $\MSE = \text{Sesgo}^2 + \text{Varianza} + \sigma^2$. La complejidad baja el sesgo y sube la varianza: el óptimo es interior (la curva U).
  \item Bajar por la diagonal flexibilidad--interpretabilidad compra precisión y pierde lectura; las herramientas de la sesión 7 recuperan parte de la lectura.
\end{enumerate}
\end{frame}

%%=================================================
\begin{frame}
\vspace{2.0cm}
\Large{Preparación para la próxima clase\ldots}
\end{frame}

%%#################################################
%%#################################################
%% PARTE 7: PRÓXIMA CLASE
%%#################################################
%%#################################################

%%=================================================
\begin{frame}{Para aprovechar al máximo la próxima sesión}
\small
\textbf{Lecturas obligatorias de esta semana:}
\begin{itemize}\itemsep4pt
  \item ISL, caps.\ 1 y 2. \\ \textcolor{gray}{Guía: la sección 2.2 es la de hoy; las figuras 2.7 y 2.12 son las de clase.}
  \item Mullainathan \& Spiess (2017), \emph{Machine Learning: An Applied Econometric Approach}, JEP. \\ \textcolor{gray}{Guía: concéntrense en la distinción $\hat{y}$ vs.\ $\hat{\beta}$ y en el ejemplo de precios de vivienda.}
\end{itemize}

\textbf{Para la primera presentación estudiantil (sesión 2):} Kim \& Zilinsky (2024), \emph{Division Does Not Imply Predictability}, \emph{Political Behavior}. \textcolor{gray}{Guía: la Figura 2 y la sección ``Growing Marginal Effects Do Not Imply Rising Joint Accuracy''.}
\pause

\vspace{0.2cm}
\textbf{Nivelación en R:} completar los módulos de introducción y \texttt{dplyr} (enlaces en la diapositiva de nivelación) \textbf{antes} de la sesión 2.
\pause

\vspace{0.2cm}
\textbf{Próxima sesión:} la regresión lineal como máquina de predecir (ISL cap.\ 3), un poco de flexibilidad (polinomios, \emph{splines}, $k$-NN; ISL cap.\ 7) y, sobre todo, cómo evaluar una predicción de regresión: RMSE, MAE, $R^2$ fuera de muestra, predicho vs.\ observado y dónde se concentra el error. \textbf{Aplicación en R:} precios de vivienda en Bogotá.
\end{frame}

%%=================================================
%% Referencias
\begin{frame}{Referencias esenciales}
\footnotesize
\begin{itemize}\itemsep2pt
  \item James, G., Witten, D., Hastie, T., \& Tibshirani, R. (2021). \emph{An Introduction to Statistical Learning} (2.\ ed.). Springer. Caps.\ 1--2. [ISL]
  \item Breiman, L. (2001). Statistical Modeling: The Two Cultures. \emph{Statistical Science}, 16(3), 199--231.
  \item Shmueli, G. (2010). To Explain or to Predict? \emph{Statistical Science}, 25(3), 289--310.
  \item Varian, H. (2014). Big Data: New Tricks for Econometrics. \emph{JEP}, 28(2), 3--28.
  \item Mullainathan, S., \& Spiess, J. (2017). Machine Learning: An Applied Econometric Approach. \emph{JEP}, 31(2), 87--106.
  \item Kleinberg, J., Ludwig, J., Mullainathan, S., \& Obermeyer, Z. (2015). Prediction Policy Problems. \emph{AER P\&P}, 105(5), 491--495.
  \item Kleinberg, J., Lakkaraju, H., Leskovec, J., Ludwig, J., \& Mullainathan, S. (2018). Human Decisions and Machine Predictions. \emph{QJE}, 133(1), 237--293.
  \item Kim, S.-y.\ S., \& Zilinsky, J. (2024). Division Does Not Imply Predictability. \emph{Political Behavior}, 46, 67--87.
  \item Gelvez, J.\ D., Cardiles, A., Martínez-González, E.\ F., \& Muñoz, M. (2026). How Predictable Is an Election? A Machine-Learning Approach to Electoral Behavior in Colombia. Documento de trabajo.
  \item Ginsberg, J., et al.\ (2009). Detecting Influenza Epidemics Using Search Engine Query Data. \emph{Nature}, 457, 1012--1014.
  \item Lazer, D., Kennedy, R., King, G., \& Vespignani, A. (2014). The Parable of Google Flu. \emph{Science}, 343(6176), 1203--1205.
\end{itemize}
\normalsize
\end{frame}


%%#################################################
%%#################################################
%% APÉNDICE: PARA PROFUNDIZAR
%%#################################################
%%#################################################
\appendix
\section{Apéndice: para profundizar}
\begin{frame}[noframenumbering]
\tableofcontents[currentsection]
\end{frame}

%%=================================================
\begin{frame}{A1. La media condicional minimiza el riesgo cuadrático}
\small
Condicionemos en $X = x$ y busquemos la mejor predicción constante $c$:
\vspace{0.1cm}

\textbf{1. Expandimos:}
\[ \E\big[(Y - c)^2 \mid X = x\big] = \Var(Y \mid X = x) + \big(\E[Y \mid X = x] - c\big)^2 \]
\pause
\textbf{2. Minimizamos en $c$:} el primer término no depende de $c$; el segundo es $\geq 0$ y se anula en
\[ c^\ast \;=\; \E[Y \mid X = x] \]
\pause
\textbf{3. Conclusión:}
\[ f^\ast(x) \;=\; \E[Y \mid X = x] \;=\; \argmin_{f} R(f) \]
\pause
\begin{block}{Interpretación}
La \textbf{media condicional} es el objetivo teórico de todo el curso. Por eso la regresión --- que la aproxima --- es el punto de partida natural del ML.
\end{block}
\end{frame}

%%=================================================
\begin{frame}{A2. Error reducible e irreducible: la cuenta}
\small
Tomemos cualquier predictor $\hat{f}$ (fijo) y evaluemos en $X = x$ (fijo), con $Y = f(x) + \varepsilon$:
\vspace{0.1cm}

\textbf{1. Sustituimos y expandimos:}
\[ \E\big[(Y - \hat{f}(x))^2\big] = \E\big[(f(x) - \hat{f}(x) + \varepsilon)^2\big] = (f(x) - \hat{f}(x))^2 + 2(f(x) - \hat{f}(x))\,\E[\varepsilon] + \E[\varepsilon^2] \]
\pause
\textbf{2. Como $\E[\varepsilon] = 0$, el término cruzado desaparece:}
\[ \E\big[(Y - \hat{f}(x))^2\big] \;=\; \underbrace{\big(f(x) - \hat{f}(x)\big)^2}_{\text{\textbf{reducible}}} \;+\; \underbrace{\Var(\varepsilon)}_{\text{\textbf{irreducible}}} \]
\pause

\begin{itemize}\itemsep4pt
  \item \textbf{Reducible:} qué tan lejos está nuestra $\hat{f}$ de la verdadera $f$. Es lo que un mejor método puede mejorar.
  \item \textbf{Irreducible} ($\sigma^2$): ni el modelo perfecto ($\hat{f} = f$) lo elimina. Es el techo de lo aprendible con estos predictores.
\end{itemize}
\end{frame}

%%=================================================
\begin{frame}{A3. Derivación sesgo--varianza (I): separar el error irreducible}
\small
\textbf{Punto de partida} ($\varepsilon_0$ es independiente de $\mathcal{D}$, con $\E[\varepsilon_0] = 0$, $\Var(\varepsilon_0) = \sigma^2$):
\[ \MSE(x_0) \;=\; \E\big[(Y_0 - \hat{f}(x_0))^2\big], \qquad Y_0 = f(x_0) + \varepsilon_0 \]
\pause

\textbf{1. Sustituimos $Y_0$ y expandimos el cuadrado:}
\[ \MSE(x_0) = \E\big[(f(x_0) - \hat{f}(x_0))^2\big] \;+\; 2\,\E\big[(f(x_0) - \hat{f}(x_0))\,\varepsilon_0\big] \;+\; \E[\varepsilon_0^2] \]
\pause

\textbf{2. El término cruzado se anula:} $\varepsilon_0$ es un ruido \emph{nuevo}, independiente de la muestra con que se construyó $\hat{f}$, y tiene media cero:
\[ \E\big[(f(x_0) - \hat{f}(x_0))\,\varepsilon_0\big] = \E\big[f(x_0) - \hat{f}(x_0)\big] \cdot \underbrace{\E[\varepsilon_0]}_{=\,0} = 0 \]
\pause

\textbf{3. Resultado parcial:}
\[ \MSE(x_0) \;=\; \E\big[(f(x_0) - \hat{f}(x_0))^2\big] \;+\; \sigma^2 \]

El primer término es el \textbf{error de estimación}; nos falta descomponerlo.
\end{frame}

%%=================================================
\begin{frame}{A4. Derivación sesgo--varianza (II): sesgo al cuadrado más varianza}
\small
\textbf{4. Sumamos y restamos $\E[\hat{f}(x_0)]$} dentro del error de estimación:
\begin{align*}
\E\big[(f - \hat{f})^2\big]
&= \E\Big[\big( \underbrace{f - \E[\hat{f}]}_{\text{constante}} + \underbrace{\E[\hat{f}] - \hat{f}}_{\text{aleatorio}} \big)^2\Big] \\
\onslide<2->{&= \big(f - \E[\hat{f}]\big)^2 + \E\Big[\big(\hat{f} - \E[\hat{f}]\big)^2\Big] + 2\big(f - \E[\hat{f}]\big)\underbrace{\E\big[\E[\hat{f}] - \hat{f}\big]}_{=\,0}}
\end{align*}
\onslide<3->{
\textbf{5. Juntando todo:}
\[ \MSE(x_0) \;=\; \underbrace{\big(f(x_0) - \E[\hat{f}(x_0)]\big)^2}_{\text{\textbf{Sesgo}}^2} \;+\; \underbrace{\Var\big(\hat{f}(x_0)\big)}_{\text{\textbf{Varianza}}} \;+\; \underbrace{\sigma^2}_{\text{\textbf{Irreducible}}} \]
}
\onslide<4->{
\begin{block}{Las tres piezas del error de predicción}
\textbf{Sesgo:} error sistemático del método (rigidez). \quad \textbf{Varianza:} sensibilidad a la muestra particular. \quad $\sigma^2$: ruido que nadie puede eliminar.
\end{block}
}
\end{frame}

%%=================================================
%% End document
\end{document}
